\chapter*{Список сокращений}             % Заголовок
\addcontentsline{toc}{chapter}{Список сокращений}  % Добавляем его в оглавление
\printnomenclature[3.5cm] % Значение ширины столбца с обозначениями стоит подбирать вручную

\begin{description}[style=multiline,leftmargin=4.7cm,labelwidth=4.4cm,labelsep=0.3cm]
    \item[\mbox{АМР}] --- автономный мобильный робот;
    \item[\mbox{MEC}] --- периферийные вычисления с множественным доступом;
    \item[\mbox{ЦП}] --- центральный процессор;
    \item[\mbox{ГП}] --- графический процессор;
    \item[\mbox{SLAM}] --- одновременная локализация и построение карты;
    \item[\mbox{V2X}] --- взаимодействие транспортного средства с окружением;
    \item[\mbox{VEC}] --- периферийные вычисления в транспортных сетях;
    \item[\mbox{VCG}] --- механизм Викри--Кларка--Гровса;
    \item[\mbox{DSIC}] --- доминантно-стратегическая совместимость;
    \item[\mbox{CTDE}] --- централизованное обучение с децентрализованным исполнением;
    \item[\mbox{RL}] --- обучение с подкреплением;
    \item[\mbox{MARL}] --- многоагентное обучение с подкреплением;
    \item[\mbox{GRU}] --- вентильный рекуррентный блок;
    \item[\mbox{DQN}] --- глубокая сеть оценки действий;
    \item[\mbox{MDP}] --- марковский процесс принятия решений;
    \item[\mbox{POMDP}] --- частично наблюдаемый марковский процесс принятия решений;
    \item[\mbox{Dec-POMDP}] --- децентрализованный частично наблюдаемый марковский процесс принятия решений;
    \item[\mbox{VDN}] --- сети разложения ценности;
    \item[\mbox{COMA}] --- контрфактический многоагентный градиент стратегии;
    \item[\mbox{MADDPG}] --- многоагентный глубокий детерминированный градиент стратегии;
    \item[\mbox{MAPPO}] --- многоагентная оптимизация стратегии с ограниченным обновлением;
    \item[\mbox{QMIX}] --- монотонное смешивание функций ценности;
    \item[\mbox{QPLEX}] --- выразительная факторизация совместной функции ценности;
    \item[\mbox{AMPLE-AMR}] --- предлагаемый гибридный метод адаптивного распределения вычислительных задач автономных мобильных роботов;
    \item[\mbox{C-AMPLE-AMR}] --- кластерный вариант метода AMPLE-AMR.
\end{description}
